\bookchapter{Anhang H: Versions-Roadmap}

\chapterintro{Dieses Buch wird laufend weiterentwickelt. Hier stehen die geplanten Richtungen für künftige Versionen – sag uns gern per Feedback, welche Inhalte du dir am meisten wünschst.}

\textbf{v1.8 ist umgesetzt: neu sind Teil 7, 8 und 9 – „KI-Programmierung in der Praxis“, „Android-Entwicklung in der Praxis“ und „iOS-Entwicklung in der Praxis“ (Kapitel 37–69, insgesamt 33 Kapitel)}\par

Von der Paradigmenrevolution der KI-Programmierung über die fünf Denkprinzipien, Kontext- und Requirements-Engineering, das Panorama von Tools und Modellen 2026, die API-Praxis und das Konversations-Engineering bis zu drei vollständigen Projekten, Qualitätssicherung und dem Weg nach vorn – so steigst du vom „KI-Nutzer“ zum „Regisseur, der die KI Code schreiben lässt“; dazu die komplette Entwicklungskette auf Android und iOS: Umgebungsaufbau (Android Studio / Xcode), Sprach-Crashkurs (Kotlin / Swift), UI-Frameworks (Jetpack Compose / SwiftUI), Daten und Netzwerk, Praxisprojekte (Haushaltsbuch + Wetterabfrage), Testen und Veröffentlichen (Google Play / App Store) sowie ein Spickzettel – so wächst du vom digitalen Nutzer ganz bis zum App-Entwickler.

\begin{tcbinfo}{Geplante Richtungen für die nächste Version (v1.9)}
\begin{itemize}
  \item Handy und Wearables: der Leitfaden für die enge Zusammenarbeit von Apple Watch, Huawei-Uhren und Smart Rings (Oura/Galaxy Ring)
  \item Reisen mit dem Handy: eSIM in der Praxis, internationale Datentarife, Offline-Karten, Bezahlen im Ausland, Vergleich von Übersetzungs-Apps
  \item Handy und Gesundheitsmanagement: Sport und Fitness, Schlafüberwachung, Herzfrequenz und Sauerstoffsättigung, Frauengesundheit, Apps für die mentale Gesundheit
  \item Kurzvideos mit KI: Praxis-Tutorials zu CapCut AI, Kling AI, digitalen Avataren und KI-Voiceover
  \item Einstieg in vertikale Kurzserien: Drehbuch-Aufbau, Aufnahme-Tricks, Plattformregeln
\end{itemize}

\end{tcbinfo}

\begin{tcbwarn}{Hinweis zur Aktualität}
Die Handy-Technik entwickelt sich extrem schnell; der Inhalt dieses Buches basiert auf dem Technologie- und Marktstand zur Jahresmitte 2026. KI-Fähigkeiten, Chipmodelle, Systemversionen und Preisspannen können sich bis zu deiner Lektüre bereits geändert haben. Bei konkreten Bedienungsschritten zählt die tatsächliche Systemversion deines Geräts; bei Kaufentscheidungen schau bitte in aktuelle Tests und Nutzerstimmen.

\end{tcbwarn}
